% mean_tau_index.tex -- round 440, the averaged tau-winding
% as an exact integer object (reviewer DCCCIV fallback lane).
% IDENTITIES_EXACT / COLLAR_BOSS_QUANTIFIED /
% BLOCK_MEAN_CENSUS / UNCONDITIONAL_OPEN.
% Builds with pdflatex.
% Companion machine check: verify_mean_tau.py.
% Discovery: experiments/tfpt-discovery/mean_tau_index_probe.py.
% NO RH CLAIM.  Research documentation, not a theorem of RH.
\documentclass[11pt]{article}

\usepackage[a4paper,margin=2.7cm]{geometry}
\usepackage{amsmath,amssymb,amsthm}
\usepackage{microtype}
\usepackage{booktabs}
\usepackage{xcolor}
\usepackage[colorlinks=true,linkcolor=blue!50!black,urlcolor=blue!50!black]{hyperref}

\newtheorem{lemma}{Lemma}
\newtheorem{prop}{Proposition}
\theoremstyle{remark}
\newtheorem{remark}{Remark}

\newcommand{\Rd}{R^{\dagger}}
\newcommand{\Td}{\mathfrak{T}^{\dagger}}
\newcommand{\Gd}{G^{\dagger}}
\newcommand{\taud}{\tau^{\dagger}}
\newcommand{\kapd}{\kappa^{\dagger}}
\newcommand{\indm}{\operatorname{ind}_{-}}
\newcommand{\half}{\tfrac12}

\title{\bfseries The mean tau index\\[2mm]
\large $\kapd$ as an integer winding; identities first,\\
no unconditional block-mean bound}
\author{Stefan Hamann}
\date{August 30, 2026}

\begin{document}
\maketitle

\begin{abstract}\noindent
Write $\Gd=(\Td)^{*}\Td\succeq 0$ and
$\taud(s)=\det(I-s\Gd)$.  All zeros of $\taud$ are real and
positive, at $s=1/\lambda_{j}(\Gd)$.
Three finite identities hold, with no analytic majorant:
(T1)~$\kapd:=\indm(\Rd-\half I)=\#\{s\in(0,1):\taud(s)=0\}$,
and $\Rd\succeq\half I$ if and only if $\kapd=0$;
(T2)~the argument principle on a contour about $(0,1)$;
(MI2)~the block mean of $\kapd$ is the winding of the averaged
integrand.
The Lean landing site
\texttt{exists\_index\_zero\_of\_block\_mean\_lt\_one} (r430)
then turns $\liminf$ of a block mean $<1$ into an index-zero
window.
The honest obstruction to (T2) as a numerical contour is
quantified: the $1/2$-cluster of $\Rd$ maps to zeros of $\taud$
in a shrinking collar of $s=1$.
No RH claim.
\end{abstract}

\medskip\noindent
\textbf{Status.}\enspace
Lemmas~\ref{lem:t1}--\ref{lem:mi2} and~\ref{lem:src} are
\textbf{SATZ} (finite algebra over $\mathbb{Q}$, machine
precision on the named windows).
Proposition~\ref{prop:collar} is a measured anatomy of the
contour-approach problem, not a bound.
Proposition~\ref{prop:census} is a sealed census.
The unconditional block-mean bound is \emph{not} claimed
(Reviewer-R439).
Ausgang
\texttt{IDENTITIES\_EXACT / COLLAR\_BOSS\_QUANTIFIED / BLOCK\_MEAN\_CENSUS / UNCONDITIONAL\_OPEN}.
No $L^{*}$ claim.  No $R^{\dagger}$ claim.  No RH claim.

\section{T1: inertia is a zero count}\label{sec:t1}

The r409 dictionary is $\Rd=(I+\Gd)^{-1}$ with
$\Gd=(\Td)^{*}\Td\succeq 0$.  Spectral calculus gives
\[
\Rd-\half I
=\frac{I-\Gd}{2(I+\Gd)}.
\]
The factor $2(I+\Gd)\succ 0$, so
$\indm(\Rd-\half I)=\#\{\lambda(\Gd)>1\}$.
The zeros of $\taud(s)=\det(I-s\Gd)$ sit at $s=1/\lambda_{j}$,
hence in $(0,1)$ precisely when $\lambda_{j}>1$.

\begin{lemma}[T1]\label{lem:t1}
$\kapd=\indm(\Rd-\half I)=\#\{s\in(0,1):\taud(s)=0\}$.
In particular $\Rd\succeq\half I$ if and only if $\kapd=0$.
\end{lemma}

\begin{proof}
The displayed congruence, plus $\Gd\succeq 0$.
Over $\mathbb{Q}$: $G=\mathrm{diag}(2,1/2)$ gives
$R=\mathrm{diag}(1/3,2/3)$, $R-\half I=\mathrm{diag}(-1/6,1/6)$,
$\indm=1$, and $\tau(s)=(1-2s)(1-s/2)$ has zeros $\{1/2,2\}$.
The r409 five-atom toy has
$C=\bigl(\begin{smallmatrix}97&23\\23&125/16\end{smallmatrix}\bigr)$,
$\det C=3661/16$, both eigenvalues of $C$ strictly above $1$,
hence $\kappa=0$ and $G=C^{-1}$ at residual $1.2\cdot 10^{-16}$.
On w9 the three counters (eigenvalues of $\Gd$ above $1$,
zeros of $\taud$ in $(0,1)$, inertia of $\Rd-\half I$) agree
at $0$, with dictionary residuals $2\cdot 10^{-14}$.
\end{proof}

Unaugmented $T_{0}$ on w9 still has $\mathrm{nneg}(I-T_{0}^{*}T_{0})=1$
(the P1 square).  The dagger repairs it: $\|\Td\|=0.99991709<1$.

\section{T2: argument principle, and the collar}\label{sec:t2}

Jacobi/Fredholm variation is the integrand of (T2):
\[
\partial_{s}\log\taud(s)
=-\operatorname{tr}\bigl(\Gd(I-s\Gd)^{-1}\bigr).
\]

\begin{lemma}[T2]\label{lem:t2}
For a simple closed contour $\Gamma$ about $(0,1)$ that
encloses no other zeros of $\taud$,
$\kapd=(1/2\pi i)\oint_{\Gamma}\partial_{s}\log\taud\,ds$.
\end{lemma}

\begin{proof}
Argument principle for a polynomial with only real positive
zeros.  The integrand identity is the Jacobi formula; a
centred finite difference at $s=0.3$ on w9 matches to relative
$4.8\cdot 10^{-8}$.
On the $\mathbb{Q}$ $2\times 2$, a circle of centre $1/2$ and
radius $0.49$ winds $1$ to $10^{-15}$.
\end{proof}

The numerical contour on a window is a different question.
Eigenvalues of $\Rd$ in $[1/2,3/5]$ map to zeros of $\taud$
near $s=1$ via $s=r/(1-r)$.  A circle of radius $0.40$ about
$1/2$ (the interval $(0.10,0.90)$) winds $0$ on every living
window measured: \emph{none} of $\kapd$ lives away from $s=1$.
Pushing the radius to $0.499$ produces a fractional winding
(w9: $-1.77$), because the vertical side of any rectangle
about $(0,1)$ passes within $O(\mathrm{gap}_{+})$ of the first
exterior zero.  For a living window
$\mathrm{gap}_{+}=1/\lambda_{\max}-1$; for a dead window the
interior zero sits at $1-\mathrm{gap}_{-}$ with
$\mathrm{gap}_{-}$ of the same order.

\begin{prop}[Collar anatomy]\label{prop:collar}
On the selected sequence $k=3,4,5,6,7,9$,
$\mathrm{gap}_{+}$ falls from $2.98\cdot 10^{-3}$ to
$1.91\cdot 10^{-8}$.
On core-42, $\mathrm{gap}_{+}\in[1.42\cdot 10^{-7},1.66\cdot 10^{-4}]$
and $n_{\mathrm{half}}:=\#\{\lvert\lambda(\Rd)-1/2\rvert\le 0.05\}$
runs $21$ to $173$.
Dead $\chi$ place their one interior zero in
$\mathrm{gap}_{-}\in[4.8\cdot 10^{-6},1.0\cdot 10^{-3}]$.
Any contour that would count $\kapd$ must thread this
shrinking corridor; that is the boss problem of (T2), not a
failure of the identity.
\end{prop}

\section{Source form of the integrand}\label{sec:src}

The IIKS/Uvarov layer supplies the characteristic identity
$\det(I-sQ)=\det(I-sE)$ and the Uvarov ratio for $\tau$-quotients.
Here $Q=\Gd$ and the source-side Gram $H=\Td(\Td)^{*}$ (a kernel
on the $X$-nodes) has the same non-zero spectrum.

\begin{lemma}[Source integrand]\label{lem:src}
$\partial_{s}\log\taud(s)=-\operatorname{tr}\bigl(H(I-sH)^{-1}\bigr)$,
equivalently the Dirichlet diagonal
$-\sum_{x\in X}[H(I-sH)^{-1}]_{xx}$.
On the unaugmented interpolant,
$\partial_{s}\log\tau_{0}(s)=-\operatorname{tr}\bigl((K_{d_{0}}[Y]W_{Y}-sI)^{-1}\bigr)$,
where $K_{d_{0}}[Y]$ is the Christoffel--Darboux kernel of the
$X$-measure at the hole nodes.
\end{lemma}

\begin{proof}
Cyclic property of the trace, plus
$\det(I_{Y}-s(\Td)^{*}\Td)=\det(I_{X}-s\Td(\Td)^{*})$.
The CD form is $G=D^{-1}K_{Y}^{-1}D^{-1}$ with $D=W_{Y}^{1/2}$,
checked on w9 at residual $2\cdot 10^{-12}$; the two traces
agree at $2\cdot 10^{-11}$.
Cyclic and slogdet identities for $\Td$ hold at $1.4\cdot 10^{-14}$.
\end{proof}

This is the raw material for a later mean-value theorem on
Dirichlet polynomials in the window data.  It is not itself a
bound.

\section{MI2 and the block-mean census}\label{sec:mi2}

\begin{lemma}[MI2]\label{lem:mi2}
On any common contour,
\[
\frac1K\sum_{k=1}^{K}\kapd_{k}
=\frac1{2\pi i}\oint\partial_{s}\Bigl(\frac1K\sum_{k=1}^{K}\log\taud_{k}\Bigr)\,ds.
\]
\end{lemma}

\begin{proof}
Linearity of the integral.  On the $\mathbb{Q}$ pair
$\mathrm{diag}(2,1/2)$ and $\mathrm{diag}(0.4,0.3)$, a common
circle of radius $0.40$ gives mean winding $1/2$ equal to the
averaged-integrand winding.  On $\{w9,\mathrm{kz}\,5\}$ the
same identity holds at $4\cdot 10^{-15}$.
The Lean arithmetic
\texttt{exists\_index\_zero\_of\_block\_mean\_lt\_one}
then yields an index-zero window in any block whose mean is
strictly less than $1$.
\end{proof}

\begin{prop}[Census]\label{prop:census}
Selected $k=3,4,5,6,7,9$: $\kapd=0$ on $6/6$, block mean $0<1$
($k=8$ not rebuilt; r421 class).
Core-42: $\kapd=0$ on $42/42$, mean $0$; no growth of
$\kapd\ge 1$ with $k$.
Dead $\chi$ $6/6$: $\kapd=1$ exactly
($\chi_{3}$-$15/19/23/33/39$ and $\chi_{4}$-$20$).
Living $\chi_{3}$-$9$: $\kapd=0$.
\end{prop}

The selected sequence is entirely living.  The empirical mean
is far below $1$.  The honest remaining question is an
\emph{unconditional} bound that keeps the mean below $1$ on
infinitely many blocks --- not claimed here.

\section{Kills and delimiters}\label{sec:kills}

\begin{itemize}
\item \textbf{$s$-No-Go.}
The r265 no-go closed $s$ as a monotone positivity dynamics
(the endpoint \emph{was} the budget).  Here $s$ only counts
zeros.  Witness: $\taud(1/2)>0$ on both living w9 and dead
$\chi_{3}$-$15$ (the $\kappa$-zero sits in the collar, not at
mid-interval).
\item \textbf{Soft mutants.}
$\operatorname{tr}\Gd$ is $50.9$ on w9 ($\kapd=0$) and $53.0$
on dead-15 ($\kapd=1$).  The r398 moment
$M_{2}=16\operatorname{tr}((I-\Rd)^{4})$ is $32.96$ vs $39.01$,
both $\gg 2$.  The $1/2$-cluster dominates every even moment;
moments do not see the threshold.
\item \textbf{Scramble.}
Unaugmented $\mathrm{nneg}(R)=21$ (bulk already dead).  The
dagger is not the story.
\item \textbf{Circularity gate.}
A pointwise bound $\lvert\sum\Lambda n^{-1/2-it}\rvert\ll X^{\varepsilon}$
is forbidden as a mean-value carrier (it would close the route
by assuming the conclusion at the prime-sum scale).
\end{itemize}

\section*{Claim boundary}

Nothing in this note is evidence for or against the Riemann
hypothesis, in either direction.  The identities are finite
linear algebra.  The census is a finite table.  The
unconditional block-mean bound is open.

\begin{thebibliography}{9}
\bibitem{r265} Round 265, \texttt{s\_monotonicity}: $s$-coordinate
  class closed (naive restatement; IIKS definite dynamics, wrong
  endpoint).
\bibitem{r398} Round 398, \texttt{high\_moment\_inertia}: even
  moments blinded by the $1/2$-cluster.
\bibitem{r409} Round 409, \texttt{borodin\_birkhoff\_intertwiner}:
  $\Td$ constructor and $R=(I+G)^{-1}$.
\bibitem{r430} Round 430, \texttt{RH/FrequentlySelected.lean}:
  \texttt{exists\_index\_zero\_of\_block\_mean\_lt\_one}.
\bibitem{r433} Round 433, \texttt{edge\_redheffer}: last pivot
  $\delta=1-q^{\dagger}$.
\bibitem{r439} Reviewer-R439 (not this round): unconditional
  block mean.
\end{thebibliography}

\end{document}
